\documentclass[a4paper]{article}

% Basic packages
\usepackage[fontset=fandol]{ctex}
\usepackage{amsmath, amssymb}
\usepackage{graphicx}
\usepackage[margin=2.5cm]{geometry}
\usepackage[most]{tcolorbox}
\usepackage{etoolbox}
\usepackage{listings}
\usepackage{hyperref}
\hypersetup{hidelinks}
\usepackage{booktabs}
\usepackage{subcaption}
\usepackage{float}
\usepackage{tikz}
\usetikzlibrary{positioning}
\IfFileExists{pgfplots.sty}{
    \usepackage{pgfplots}
    \pgfplotsset{compat=1.18}
}{}

% Highlight boxes
\newtcolorbox{knowledgebox}[1]{
    enhanced, breakable, colback=blue!5!white, colframe=blue!75!black, colbacktitle=blue!75!black,
    coltitle=white, fonttitle=\bfseries, title=#1, attach boxed title to top left={yshift=-2mm, xshift=2mm},
    boxrule=1pt, sharp corners
}

\newtcolorbox{importantbox}[1]{
    enhanced, breakable, colback=yellow!10!white, colframe=yellow!80!black, colbacktitle=yellow!80!black,
    coltitle=black, fonttitle=\bfseries, title=#1, sharp corners
}

\newtcolorbox{warningbox}[1]{
    enhanced, breakable, colback=red!5!white, colframe=red!75!black, colbacktitle=red!75!black,
    coltitle=white, fonttitle=\bfseries, title=#1, sharp corners
}

\newtcolorbox{dialoguebox}[1]{
    enhanced, breakable, colback=green!4!white, colframe=green!45!black, colbacktitle=green!45!black,
    coltitle=white, fonttitle=\bfseries, title=#1, sharp corners
}

% Listings
\lstset{
    language=Python,
    basicstyle=\ttfamily\small,
    keywordstyle=\color{blue},
    stringstyle=\color{red!60!black},
    commentstyle=\color{green!60!black},
    breaklines=true,
    frame=single,
    numbers=left,
    numberstyle=\tiny\color{gray},
    captionpos=b,
    extendedchars=false
}

% Front-page metadata
\newcommand{\notetitle}{Lintel 征服 UNIX：Intel、微软与 Linux\\如何拆掉垂直整合的服务器王国}
\newcommand{\noteauthors}{五道口纳什 \& Codex}
\newcommand{\notedate}{\today}
\newcommand{\videochannel}{Asianometry}
\newcommand{\videopublishdate}{2026 年 10 月 25 日}
\newcommand{\videoduration}{42 分 35 秒}
\newcommand{\videourl}{https://www.youtube.com/watch?v=ksx5g5WvAIw}
\newcommand{\videocoverpath}{source.jpg}

\begin{document}

\begin{titlepage}
\centering
\vspace{1.2cm}
{\huge\bfseries \notetitle\par}
\vspace{0.8cm}
{\large \noteauthors\par}
\vspace{0.3cm}
{\large \notedate\par}
\vspace{1.2cm}

\ifdefempty{\videocoverpath}{
{\small 请在生成笔记时填入视频封面图的本地路径。\par}
}{
\includegraphics[width=0.82\textwidth,height=0.45\textheight,keepaspectratio]{\videocoverpath}\par
}

\vfill
\begin{tcolorbox}[width=0.9\textwidth, colback=black!2!white, colframe=black!60, sharp corners]
\textbf{视频作者/频道}：\videochannel\par
\textbf{发布日期}：\videopublishdate\par
\textbf{视频时长}：\videoduration\par
\textbf{视频链接}：
\ifdefempty{\videourl}{
未填写
}{
\href{\videourl}{\nolinkurl{\videourl}}
}
\par\textbf{PDF制作Skill}：\href{https://github.com/wdkns/wdkns-skills}{\nolinkurl{https://github.com/wdkns/wdkns-skills}}
\end{tcolorbox}
\end{titlepage}

\tableofcontents
\newpage

%% --- 正文内容开始 --- %%

\section{引言：一场被忽略的产业更替}

1990 年代初，如果你想给公司买一台能干重活的服务器，答案几乎是唯一的：一台装着 UNIX 操作系统、跑在 RISC 处理器上的专有整机。Sun 的 SPARC 加 Solaris、IBM 的 POWER 加 AIX、惠普的 PA-RISC 加 HP-UX——每一套都是软硬件打包出售的封闭体系。整套系统动辄数十万美元，而且一旦买进去，就很难再换出来。

到了 2000 年代中期，同一类负载的主流选择变成了 x86 服务器加 Linux。这个转变之所以值得单独讲一遍，是因为它并不是某一次技术突破造成的。CPU 性能不是一夜之间超过 RISC 的，Linux 也不是某天突然变得完善。真正发生的事情是：三个原本各自为战的角色——Intel 的芯片、微软的 Windows NT、以及来自自由软件社区的 Linux——在同一个方向上推动了十几年，最后把一套以垂直整合为核心的商业模式挤出了主流市场。

视频作者用了很形象的称呼：\textbf{Lintel}，即 Linux + Intel。这个合成词背后的意思是，一个原本被视作“PC 玩具”的芯片架构，配上一个没有所属公司的操作系统，最终承接了原本属于昂贵 UNIX 机器的负载。

\begin{figure}[H]
\centering
\includegraphics[width=0.92\textwidth]{figures/fig-gen-timeline.pdf}
\caption{从 GNU 项目启动到 Windows Server 收入追平 UNIX 的主要节点。图中时间取自视频正文的叙述。}
\end{figure}

这份讲义按照“问题—方案—机制—证据—结论”的顺序重述视频内容，并补充了一些视频里一带而过、但对理解机制很重要的细节。阅读时建议抓住三条并行的线索：

\begin{enumerate}
\item \textbf{成本结构}：Intel 靠 PC 的高销量摊薄了芯片成本，这是它能把服务器芯片做便宜的底层原因；
\item \textbf{软件生态}：Windows NT 提供了企业级特性，但它的发布节奏和可移植性策略又处处与 Intel 的利益相左；
\item \textbf{许可证治理}：GPL 让 Linux 既能吸收大公司的贡献，又不会被某家公司切成碎片。
\end{enumerate}

之后我们会看到，UNIX 阵营输掉这场竞争的原因，主要不在技术指标，而在于商业模型。

\subsection{本章小结}
1990 年代的企业服务器市场由垂直整合的 UNIX 体系主导，价格高昂且切换成本极高。Intel、微软与 Linux 三股力量从不同方向侵蚀这套体系，最终形成了 “Lintel” 这一新的主流组合。理解这个过程的关键，是把技术演进放回成本结构和商业模式的背景里看。

\section{1990 年代初的企业服务器：垂直整合的 UNIX 王国}

\subsection{服务器的分层}

在讨论竞争之前，先要把“服务器”这个词拆开。视频里给出了一个大致的三层结构，这三层对性能、可靠性和价格的要求完全不同，也因此成为后来不同技术路线的战场。

\begin{knowledgebox}{三层服务器市场}
\begin{itemize}
\item \textbf{工作站}：性能较强的桌面机，用于图形、科学计算等重负载任务。它可以顺便在局域网里提供一两个文件服务，算是低端服务器。
\item \textbf{中端服务器}：由多颗 CPU 组成的机箱级系统，运行数据库、文件存储、复杂业务应用，通过网络向大量客户端提供算力。
\item \textbf{高端大型机}：多 CPU、64 位寻址、虚拟化、大容量内存，服务于最核心的事务处理。
\end{itemize}
这里的划分不是绝对的，但它决定了一件事：不同层级的客户对“什么算够用”的容忍度不同，而这正好给了后来者从下层切入、逐层向上打的路径。
\end{knowledgebox}

到 1990 年代中后期，这些机器还多了一个用途：托管网站。随着互联网热潮扩张，网站托管成了一个巨大的市场。虽然跑一个网站并不需要高端大型机，但当时的默认做法仍然是买 UNIX 服务器——这为后来的替代留下了空间。

\subsection{三家代表与它们的共同结构}

这一时期的企业服务器市场，被几家垂直整合的系统厂商瓜分。所谓垂直整合，是指同一家公司同时提供处理器、操作系统、编译器和整机，客户买到的是一整套封闭系统。视频点名了四家：

\begin{itemize}
\item \textbf{Sun Microsystems}：SPARC 处理器 + Solaris 操作系统；
\item \textbf{IBM}：POWER 系列 RISC 处理器 + AIX 操作系统；
\item \textbf{Hewlett-Packard}：PA-RISC 处理器 + HP-UX 操作系统；
\item \textbf{Silicon Graphics}：面向图形工作站起家的另一套专有体系。
\end{itemize}

\begin{figure}[H]
\centering
\includegraphics[width=0.66\textwidth]{figures/fig02-vendor-vertical.jpg}
\caption{惠普的 PA-RISC 处理器（左）与运行 HP-UX 的工作站（右）。芯片、操作系统、整机来自同一家供应商，这正是“垂直整合”的具体形态。\protect\footnotemark}
\end{figure}
\footnotetext{视频画面时间区间：00:02:40--00:02:45。}

这些系统的价格从数十万美元起步。更关键的是它们的封闭性：编译器、驱动、系统库、技术支持合同全部绑定在同一个供应商身上。客户一旦把业务系统搭在某个平台上，迁移就意味着重写、重测和重新培训。

视频里用一句话概括了这种关系：\textbf{buy in, get locked in}（买进去，就被锁死）。

\begin{warningbox}{锁定不是合同条款，而是技术债}
很多人把“锁定”理解成法律意义上的排他协议。实际上 UNIX 厂商的锁定主要来自技术层面：私有编译器行为、非标准的系统调用扩展、与特定硬件绑定的驱动，以及只有原厂工程师才熟悉的运维知识。这些东西叠加起来，会让一次平台迁移的隐性成本远超硬件本身的差价。
\end{warningbox}

\subsection{本章小结}
1990 年代初的服务器市场是一个典型的垂直整合市场：四家厂商各自掌握从芯片到操作系统的完整栈，靠封闭性和切换成本维持高毛利。这套结构在技术上没有明显缺陷，真正的问题在于它无法享受通用计算的高销量带来的成本下降——而这正是 Intel 的强项。

\section{Intel 的算盘与 Windows NT 的登场}

\subsection{高销量经济学：Intel 的底层飞轮}

要理解 Intel 为什么敢进服务器市场，得先理解它在 PC 市场是怎么赢的。视频把它拆成一个自我强化的循环：

\begin{figure}[H]
\centering
\begin{tikzpicture}[every node/.style={font=\small},
  box/.style={draw=blue!55!black, rounded corners, fill=blue!7, minimum width=2.9cm,
              minimum height=0.9cm, align=center, inner sep=4pt}]
\node[box] (vol)    at ( 90:3.7) {高销量\\（PC 市场规模）};
\node[box] (yield)  at ( 18:3.7) {更高良率\\（同样晶圆出更多好芯片）};
\node[box] (cost)   at (-54:3.7) {更低单位成本};
\node[box] (profit) at (-126:3.7){更高利润与现金流};
\node[box] (invest) at (162:3.7) {再投资：更先进制程};
\draw[->, thick, blue!70!black] (vol) to[bend left=12] (yield);
\draw[->, thick, blue!70!black] (yield) to[bend left=12] (cost);
\draw[->, thick, blue!70!black] (cost) to[bend left=12] (profit);
\draw[->, thick, blue!70!black] (profit) to[bend left=12] (invest);
\draw[->, thick, blue!70!black] (invest) to[bend left=12] (vol);
\end{tikzpicture}
\caption{Intel 在 PC 市场的高销量飞轮。销量不仅摊薄了固定成本，还通过良率学习曲线压低了单位成本，释放出的现金流又被投入更先进的制程。}
\end{figure}

这个循环的直接后果是：AMD、Cyrix 这类竞争者即使设计出相近的产品，也很难在成本上追上，因为它们的出货量小得多，良率学习的速度也慢得多。Intel 由此在 PC 市场建立了难以逾越的领先。

于是问题变成：这套玩法能不能搬到企业服务器市场？答案取决于能否找到一个足够强的软件伙伴——因为服务器客户买的不是芯片，而是能跑业务系统的平台。

\subsection{Windows NT：一个真正的企业级操作系统}

这个伙伴是微软。1993 年，由 Dave Cutler 领导的团队发布了 Windows NT。这里的 “NT” 官方解释是 “New Technology”。理解 NT 的关键，是把它和消费级 Windows 严格区分开。

视频里有一个很形象的对比：普通的 Windows（当时是 Windows 3.x）本质上是一个\textbf{粗暴嫁接在 MS-DOS 之上的图形外壳}；而 NT 是从零开始为企业环境重写的操作系统。

\begin{figure}[H]
\centering
\includegraphics[width=0.58\textwidth]{figures/fig03-nt-sales-9394.jpg}
\caption{IDC 统计的 Windows NT 工作站版与服务器版销量（单位：千套）。1993 年服务器版仅约 3 万套，1994 年升至 11.5 万套。\protect\footnotemark}
\end{figure}
\footnotetext{视频画面时间区间：00:04:44--00:04:51。}

两者的差别体现在几个具体的地方：

\begin{itemize}
\item 内存保护与抢占式多任务：一个用户进程崩溃不会拖垮整个系统；
\item 可移植性：内核被设计成可以在多种处理器架构上编译，而不是死绑 x86；
\item 面向企业的安全模型：用户、组、权限的概念从内核层就存在。
\end{itemize}

微软对外宣传的重点也正是这三点：更强大、更可靠、更安全。

\subsection{早期并不成功}

不过 NT 的第一个版本并没有一炮而红。视频给了几个具体数字和原因：

\begin{itemize}
\item 首发版本运行在 Intel 386 上，功能可用，但缺少竞争对手（IBM 的 OS/2、Novell 的 NetWare、各种 UNIX）已经具备的能力；
\item 对当时的硬件要求偏高：需要 16 MB 内存；
\item 1993 年 NT Server 大约只卖出 3 万套。
\end{itemize}

1994 年情况好转，微软卖出了约 11.5 万份 NT 许可证。虽然仍落后于同代的工作站版本，但增长趋势已经明确。视频在这里开了个玩笑，说微软像“终结者”或者《Air Bud》的续集——不管怎样都会一直出下去。这个比喻背后是一个真实的商业事实：微软有能力在一个尚未盈利的品类里持续投入多年，直到产品成熟。

\begin{figure}[H]
\centering
\includegraphics[width=0.58\textwidth]{figures/fig04-nt-sales-9396.jpg}
\caption{同一统计口径延长到 1996 年：服务器版销量从 1993 年的 3 万套增长到 1996 年的约 67 万套（其中 1996 年为估计值）。\protect\footnotemark}
\end{figure}
\footnotetext{视频画面时间区间：00:07:04--00:07:11。}

\subsection{本章小结}
Intel 在 PC 市场建立的“高销量—高良率—低成本—再投资”飞轮，是它进军服务器的底气。但服务器客户买的是完整平台，因此 Intel 需要一个企业级操作系统作为搭档，这个角色由微软的 Windows NT 承担。NT 在 1993 年起步艰难，但它具备消费级 Windows 所没有的内存保护、抢占式多任务和可移植性，这三点决定了它后来能撑起“Wintel”的服务器路线。

\section{Pentium Pro 与“标准高容量服务器”}

\subsection{一颗改变了性能预期的芯片}

1995 年底，Intel 推出了它认为真正能打进服务器市场的第一颗芯片：32 位的 Pentium Pro。视频给出的规格对比很直观：

\begin{table}[H]
\centering
\caption{Pentium 与 Pentium Pro 的规格对比}
\begin{tabular}{lcc}
\toprule
 & Pentium & Pentium Pro \\
\midrule
晶体管数 & 310 万 & 550 万 \\
主频 & 约 60 MHz & 约 150--200 MHz \\
发布年份 & 1993 & 1995 \\
\bottomrule
\end{tabular}
\end{table}

真正让业界震动的不是频率，而是基准测试结果。在某些整数基准上，Pentium Pro 小幅超过了当时最快的 RISC 芯片。虽然后来 MIPS R10000 这类更快的 RISC 芯片很快又追了上来，但这次测试改变了一个长期存在的认知：x86 不再只是“桌面架构”。

\begin{figure}[H]
\centering
\includegraphics[width=0.78\textwidth]{figures/fig05-spec-benchmark.jpg}
\caption{1995 年的 SPECint95 / SPECfp95 对比。Pentium Pro 在整数测试上已能与同期 RISC 芯片同场竞争，但在浮点上仍明显落后。\protect\footnotemark}
\end{figure}
\footnotetext{视频画面时间区间：00:05:25--00:05:31。}

\begin{importantbox}{性能楔子的作用是“够用”，不是“最快”}
视频作者特意指出，Pentium Pro 在浮点性能上仍然落后于最好的 RISC 芯片。这在当时并不是致命的：企业服务器上跑的多数负载（数据库事务、文件服务、Web 服务）以整数和内存访问为主。真正关键的是“性价比”——只要性能够用而价格低一个数量级，客户就会重新计算。
\end{importantbox}

\subsection{把芯片变成“标准产品”}

仅仅有芯片还不够。服务器客户买的是整机，而整机的集成工作琐碎且耗时。1995 年中，Intel 开始向合作伙伴分发所谓“标准高容量服务器”（Standard High Volume Server，SHV）的参考设计板。

这些参考设计的具体规格是：

\begin{itemize}
\item 支持 1--4 颗 Pentium Pro；
\item 1--2 GB 工作内存；
\item 数百 GB 存储；
\item Intel 宣称可处理每秒 2000 笔事务。
\end{itemize}

Hewlett-Packard、Dell、Compaq、Gateway 2000 等 OEM 可以直接采用这套设计，把产品上市时间压缩到几个月。在 1995 年 11 月的发布会上，Intel 宣布基于 Pentium Pro 的整机价格区间从 2600 美元的工作站到 2 万美元的服务器。

\begin{figure}[H]
\centering
\includegraphics[width=0.62\textwidth]{figures/fig06-shv-p6.jpg}
\caption{1995 年 5 月 InfoWorld 的报道：Intel 向 OEM 推出完整的 P6 服务器系统方案。\protect\footnotemark}
\end{figure}
\footnotetext{视频画面时间区间：00:06:09--00:06:16。}

注意这里发生的事情：Intel 把自己从“卖芯片”变成了“卖整机参考方案”。这一招绕开了传统服务器厂商的系统集成能力，让一批原本只做 PC 的 OEM 也能快速进入服务器市场。等到这些 OEM 成长起来，UNIX 厂商面对的就不再是 Intel 一家，而是一整个以 PC 供应链为基础的产业群。

\subsection{1996 年的转折}

配合 1996 年发布的 Windows NT 4.0，这套组合开始见效。IDC 的数据显示，1996 年 NT 4.0 的销售额增长了 85\%，市场份额达到 27\%。竞争依然激烈，但 Wintel 在服务器市场已经建立起了类似它在 PC 市场的根据地。

\subsection{本章小结}
Pentium Pro 用“整数性能追平 RISC、价格低一个数量级”的方式，为 x86 打开了服务器市场的大门。SHV 参考设计则把芯片变成了 OEM 可以直接组装的标准化产品，绕过了传统服务器厂商的系统集成壁垒。到 1996 年，Wintel 在服务器市场已经站稳，UNIX 厂商开始感到压力。

\section{联盟的裂缝：NSP、节奏错配与 SMP 天花板}

Wintel 是一个利益共同体，但不是利益完全一致的共同体。这一章讲的是两者之间的结构性摩擦——这些摩擦在当时看起来只是商业纠纷，事后回看却是 Intel 最终转向 Linux 的直接原因。

\subsection{NSP 事件：谁拥有操作系统之上的那一层}

1995 年，Intel 推出了一个叫“原生信号处理”（Native Signal Processing，NSP）的多媒体软件标准。它的思路很直接：把一部分音频和视频计算从专用芯片搬到 CPU 上，用软件和 API 实现。对 Intel 来说，这既能让用户感受到“新 CPU 更快”，也能扩大 x86 的应用面。

微软的反应极为强烈。在微软看来，NSP 是一种“中间件”，它绕过了操作系统直接接触硬件和应用，等于侵入了 Windows 的地盘。微软高管明确表示不会支持 NSP，比尔·盖茨还专门给安迪·格鲁夫发邮件表达不满。

\begin{figure}[H]
\centering
\includegraphics[width=0.78\textwidth]{figures/fig07-nsp.jpg}
\caption{1995 年 10 月 PC Magazine 的报道：在与微软的争执之后，Intel 大幅调整了 NSP 参考平台。\protect\footnotemark}
\end{figure}
\footnotetext{视频画面时间区间：00:07:17--00:07:23。}

事情最终以 Intel 退让收场。NSP 事件后来因为微软的反垄断诉讼而被公开，成为微软利用平台地位打压合作方的例证之一。两家公司对外都淡化处理，但这件事给 Intel 留下了一个清晰的信号：\textbf{它的“伙伴”随时有能力用平台权力否决它的战略。}

\subsection{NT 的可移植性：一把双刃剑}

第二个摩擦点更隐蔽，但影响更深远。与消费级 Windows 不同，Windows NT 从设计之初就是可移植的。虽然绝大多数 NT 实例跑在 x86 上，但它同样支持 DEC 的 Alpha、IBM 与 Apple 的 PowerPC，以及 MIPS 架构。

\begin{figure}[H]
\centering
\includegraphics[width=0.42\textwidth]{figures/fig08-nec-nt-portable.jpg}
\caption{1993 年 NEC 在 InfoWorld 上投放的广告，标题是“The brains behind Windows NT”，宣传其 MIPS R4400 芯片可以运行 NT。\protect\footnotemark}
\end{figure}
\footnotetext{视频画面时间区间：00:08:40--00:08:46。}

对客户来说，这是好事：买 NT 不会被绑在某一种芯片上。对 Intel 来说，这是一个战略隐患。微软还在不断签署新的芯片支持协议，等于在 Intel 身边布下了一圈“备胎”。于是一个问题浮出水面：

\begin{importantbox}{Intel 的核心恐惧}
如果 NT 成为服务器操作系统的事实标准，那么微软完全有能力把 x86 降格为“众多硬件架构之一”。到那时，Intel 在 PC 市场积累的规模优势，就会被稀释到一个多架构竞争的市场里。这相当于把 Intel 最重要的护城河交给合作伙伴来决定。
\end{importantbox}

\subsection{节奏错配：摩尔定律遇上三年一版的 Windows}

就算不考虑战略博弈，两家公司的产品节奏也天然不同。

Intel 遵循摩尔定律，大约每两年就能推出更快、更好的芯片。微软则大约每三到五年才发布一个新版 Windows，而且新系统未必会立刻用满新芯片的能力。

视频给了两个非常有说服力的例子：

\textbf{第一个例子：32 位。} Intel 在 1985 年就推出了 32 位的 386 处理器。微软直到 1995 年的 Windows 95 才拿出面向消费者的 32 位操作系统——整整晚了十年。

\textbf{第二个例子：Chicago 的图形层。} 在 Pentium Pro 的开发过程中，微软在多次会议上告诉 Intel，下一代 Windows（代号 Chicago，即后来的 Windows 95）将是纯 32 位、去掉全部遗留的 16 位代码。Intel 据此把 Pentium Pro 的优化重点放在了 32 位性能上。结果 1994 年微软公开宣布，Chicago 的图形显示接口——操作系统的关键部分——仍将保留 16 位。

\begin{figure}[H]
\centering
\includegraphics[width=0.8\textwidth]{figures/fig09-pentiumpro-16bit.jpg}
\caption{1995 年 11 月《Daily Record》的报道标题：“Pentium Pro sets PC standard — But Intel's speedy new chip lags on older 16-bit software”。\protect\footnotemark}
\end{figure}
\footnotetext{视频画面时间区间：00:10:36--00:10:41。}

此时 Pentium Pro 已经进入开发后期，改设计来不及了。结果芯片上市后，媒体在消费级桌面场景下测出了糟糕的图形性能，Intel 替微软的设计决策背了锅。

\subsection{SMP 天花板：4--8 路对 64 路}

第三个矛盾出现在服务器场景的核心指标上：多处理器扩展能力。

NT 4.0 单个操作系统实例最多支持 4--8 颗处理器，而当时一些 UNIX 变体可以支持到 64 颗。这意味着 NT 在架构上被挡在了高端市场之外——客户如果要在单机里跑大型数据库，仍然只能选 UNIX。

\begin{figure}[H]
\centering
\includegraphics[width=0.92\textwidth]{figures/fig-gen-smp.pdf}
\caption{单个操作系统实例支持的处理器数对比（数据取自视频正文）。NT 4.0 的上限、Linux 2.2 的上限，以及随后两代 Linux 内核的改进构成了后来竞争的关键变量。}
\end{figure}

\subsection{两种增长模型}

为什么会出现这种错配？视频给出了一个很有洞察力的解释：两家公司的增长模型不同。

\begin{table}[H]
\centering
\caption{Intel 与微软的增长模型对比}
\begin{tabular}{p{0.2\textwidth}p{0.34\textwidth}p{0.34\textwidth}}
\toprule
 & Intel & Microsoft \\
\midrule
收入来源 & 卖最新的芯片 & 服务整个已安装基础 \\
对软件的要求 & 需要能吃掉新算力的软件 & 需要向后兼容、覆盖尽可能多用户 \\
发布节奏 & 约每两年一代 & 约三到五年一版 \\
对旧平台的态度 & 希望用户尽快换新 & 必须照顾存量客户 \\
\bottomrule
\end{tabular}
\end{table}

微软并非不愿意投入——开发一个新版 Windows 需要数千名员工、可能耗资数十亿美元。问题在于软件本身的属性决定了它必须照顾存量。这两种模型各有道理，但放在一起就会持续产生摩擦。

那么 Intel 有什么选择？视频在这里点题：1990 年代后期，一个非常有力的新候选出现在了 Intel 的雷达上——Linux。

\subsection{本章小结}
Wintel 联盟内部存在三重结构性矛盾：NSP 事件表明微软会用平台权力否决 Intel 的越界；NT 的可移植性意味着微软可以把 x86 降格为普通架构；产品节奏与 SMP 上限的错配则让 Intel 的新芯片无法在服务器市场充分变现。这些矛盾共同构成了 Intel 寻找“第二操作系统”的动机。

\section{GNU 与 Linux：一个没有所有者的替代品}

在讲 Intel 如何发现 Linux 之前，需要先讲清楚 Linux 是怎么来的。视频把这段历史压缩在几分钟里，但其中的每个环节都直接影响 Linux 后来的治理结构。

\subsection{Stallman 的动机：一个被杀死的社区}

1984 年 1 月，麻省理工学院（MIT）的程序员 Richard Stallman 正式启动了一个替换 UNIX 的项目。要理解他为什么这么做，得回到 MIT 人工智能实验室的那段历史。

Lisp 最初是早期人工智能研究中的优雅工具。当商人意识到 Lisp 和它所支撑的 AI 可以成为一个大产业时，实验室分裂成了两家公司：Symbolics 和 Lisp Machines Inc.。Symbolics 融了数百万美元，试图成为“那个时代的英伟达”，把实验室里最顶尖的黑客全部挖走，只留下一个空壳。

\begin{figure}[H]
\centering
\includegraphics[width=0.72\textwidth]{figures/fig10-symbolics-lmi.jpg}
\caption{Symbolics 与 Lisp Machines Inc.（LMI）的标志，这两家公司由 MIT AI 实验室分裂而来。\protect\footnotemark}
\end{figure}
\footnotetext{视频画面时间区间：00:12:27--00:12:33。}

Stallman 始终没有原谅这件事。他后来说的一句话被视频引用：“Everything's dying.”（一切都在死去。）出于对黑客社区消亡的失落，他决定建立一个新的社区，而它的基础是“自由软件”这个概念。

\subsection{GNU：零件齐全，只缺内核}

1983 年底，Stallman 启动了一个构建自由操作系统的项目，它遵循 UNIX 的设计原则——模块化、可移植——但不复制 UNIX 的代码，以避免商业秘密问题。这个项目叫 GNU，递归缩写为 “GNU's Not UNIX”。

接下来的几年里，Stallman 和社区把工具链的各个部分都做了出来：编辑器、编译器、shell、各种工具程序。整个系统只缺一样东西：\textbf{内核}。

\begin{figure}[H]
\centering
\includegraphics[width=0.66\textwidth]{figures/fig11-kernel-diagram.jpg}
\caption{内核在系统中的位置：它位于应用与硬件之间，负责内存管理、CPU 时间分配、磁盘与设备访问。视频在此处特别强调，内核是操作系统真正的心脏。\protect\footnotemark}
\end{figure}
\footnotetext{视频画面时间区间：00:13:33--00:13:38。}

视频在这里特意解释了一下内核的重要性，因为这是理解后面所有事情的前提。内核负责内存怎么分配、每个进程能分到多少 CPU 时间、磁盘空间怎么管理。它把操作系统的其余部分与硬件隔开。缺了它，GNU 只是一套漂亮但没有心脏的工具。

\subsection{Linus 与集市模式}

内核最终由一个 21 岁的芬兰本科生补上了。Linus Torvalds 受到教学用 UNIX 系统 MINIX 的启发，在 1991 年 8 月于 MINIX 的 Usenet 新闻组发帖，说自己正在做一个自由的操作系统，并征求功能建议。帖子吸引了 MINIX 生态的贡献者，最终演化成了自己的项目：Linux。

\begin{figure}[H]
\centering
\includegraphics[width=0.78\textwidth]{figures/fig12-gpl-defining.jpg}
\caption{2021 年 CIO 的采访标题：“Linus Torvalds says GPL was defining factor in Linux's success”。Torvalds 在采访中反复强调，GPL 防止了碎片化。\protect\footnotemark}
\end{figure}
\footnotetext{视频画面时间区间：00:15:13--00:15:19。}

Linus 的角色很快从写代码转变为审阅别人的贡献并合并进代码库。这种被后来概括为“集市”的模式，让 Linux 能以比公司更快的速度交付新代码。1991 年 10 月，第一个可用版本发布；1992 年，0.12 版采用了 Stallman 的 GPL 许可证。

\subsection{GPL 的制度设计}

GPL 的具体条款值得单独说明，因为它解释了 Linux 为什么没有重蹈 UNIX 的覆辙。

GPL 要求：任何人都可以阅读和修改代码，但只要分发修改后的版本，就必须同样以开源方式发布。用更直白的话说：

\begin{importantbox}{GPL 解决了两个具体问题}
\textbf{防止碎片化}：公司不能把 Linux 内核拿过去，加点自己的东西，再作为闭源的竞争产品卖出去。这避免了 UNIX 当年分裂成三十多个互不兼容变体的命运。\\
\textbf{保证贡献回流}：任何公司对内核做的改进都会惠及所有人，不会被人锁进私有分支。这让大公司愿意投入——IBM 后来有数百名员工为内核贡献代码。\\
Torvalds 后来多次表示，选择 GPL 是他做过的最正确的决定。
\end{importantbox}

这里有一个容易被忽略的因果关系：Linux 之所以能被大公司接纳，是因为 GPL 让它们确信自己的投入不会被竞争对手独占。许可证在这里起到了产业治理工具的作用。

\subsection{本章小结}
GNU 提供了完整的工具链但缺少内核，Linux 补上了这块拼图，而 GPL 则从制度上保证了内核不会被碎片化、贡献不会被私有化。三者结合，产生了一个不属于任何公司、却能被所有公司使用的操作系统。这个治理结构是 Linux 后来能承接企业级负载的根本原因。

\section{发行版与 Red Hat：把内核变成产品}

\subsection{内核是引擎，发行版是整车}

Linux 只是一个内核。视频用了一个很贴切的比喻：内核是汽车引擎，非常重要，但你还得需要整辆车的其余部分。

对早期极客来说，自己组装系统是可以接受的。但让普通用户先“造一辆车”再开，显然不现实。于是出现了\textbf{发行版}（distribution，简称 distro）：把 Linux 内核与 GNU 项目或其他自由软件的实用工具打包在一起，形成一套稍微方便一些的操作系统套件。

视频还打了个比方：这就像高中时代大家互相传递的混音带（mixtape）——把喜欢的歌按某种顺序组合起来。好的发行版能流行到可以变现，其商业模式被形容为“卖瓶装水”：包装不同，水是同一份。

\subsection{Red Hat 的起点}

早期有代表性的发行版包括 Slackware、Debian，以及后来最重要的 Red Hat。

Red Hat 由 Marc Ewing 在 1993 年创立。Ewing 很早就开始用 Linux，但用了一段时间后发现，自己花在管理 Linux 机器上的时间比写代码还多。于是他做了一个发行版，放在他公寓的备用卧室里——那间公寓靠近杜克大学。名字来自他在大学时戴的一顶红白条纹长曲棍球帽。

Red Hat 做的事情，是把 Linux 变成一个更专业的产品。最典型的例子是 Red Hat Package Manager（RPM）：它让用户可以方便地安装、删除、更新软件包，比如 Web 服务或 FTP 服务。

\begin{figure}[H]
\centering
\includegraphics[width=0.46\textwidth]{figures/fig13-redhat-box.jpg}
\caption{Red Hat Linux 6.1 的零售包装盒。发行版可以刻成光盘、装盒、在零售店销售——这是 Linux 走向主流用户的具体形态。\protect\footnotemark}
\end{figure}
\footnotetext{视频画面时间区间：00:18:31--00:18:36。}

1995 年，Ewing 和同事 Robert Young 成立了 Red Hat 公司，收入来自销售光盘，以及配套的培训、教育和客户支持服务。

\subsection{把 Linux 带出极客圈}

在最初的几年，Linux 的用户群体是技术水平很高的黑客和穷学生。Red Hat 以及其他商业发行版（如 Caldera）的崛起，把这个系统带出了那个小圈子，进入了企业主流视野。

这个转变的意义在于：企业客户买的从来不是“内核”，而是“一套有支持合同的、可运维的系统”。发行版公司恰好填补了这个空缺。

\subsection{本章小结}
内核需要有人把它包装成可安装、可维护、可支持的产品，发行版承担了这个角色。Red Hat 通过 RPM 包管理、光盘零售和技术支持服务，把 Linux 从黑客的玩具变成了企业可以采购的商品。这一步是 Linux 进入商业世界的前提。

\section{Linux 的早期市场：价格只是入场券}

\subsection{用户增长曲线}

Linux 的采用起步很小，但增长很快。视频给出了几个时间点的估计：

\begin{itemize}
\item 1992 年初：可能只有几百名用户；
\item 1996 年：大约 100 万到 200 万；
\item 1997 年中：全球用户保守估计在 300 万到 500 万之间。
\end{itemize}

相比之下，同期 Windows 约有 1 亿用户，Macintosh 有 5000 万以上。Linux 的绝对规模仍然很小，但增长速度完全不同。

\subsection{价格差距有多大}

价格当然是因素之一，而且差距极大。

\begin{figure}[H]
\centering
\includegraphics[width=0.9\textwidth]{figures/fig-gen-price-ladder.pdf}
\caption{同一台服务器上操作系统的价格对比（数据取自视频正文，对数刻度）。}
\end{figure}

\begin{itemize}
\item 一张 Red Hat 光盘加包装盒在零售店卖约 50 美元；
\item Windows NT 许可证从单份升级的约 150 美元，到 10 个座位的约 1100 美元不等；
\item 包含软硬件的 UNIX 系统则要 2.5 万到 25 万美元甚至更多。
\end{itemize}

视频引用了一个具体案例：1997 年《Byte》杂志报道，伊士曼柯达的一位系统顾问在某个 Photo CD 相关项目中安装了 Linux，因为不必购买高价 UNIX 替代方案，公司省下了 1 万美元。

\subsection{不只是便宜}

但真正有意思的是，企业和机构采用 Linux 并不完全是因为便宜。视频列出了几个更实质的原因：

\textbf{可靠性。} 一位用户形容 Linux 几乎“不会崩溃”（crash-proof），这直接转化为更高的生产力。

\textbf{源代码可访问。} 对高级用户来说，能读到源代码是 Windows 永远给不了的。这意味着可以按需改造系统。有人把 Linux 移植到 Palm Pilot 上；洛斯阿拉莫斯的研究人员用 Linux 把 68 台 PC 联网，用来模拟原子冲击波。

\begin{figure}[H]
\centering
\includegraphics[width=0.72\textwidth]{figures/fig14-linux-price-clip.jpg}
\caption{1996 年 5 月《Philadelphia Inquirer》的报道，讨论 Linux 的低成本与可获得性。文中提到“最好的 Linux 来源是一家叫 Red Hat 的公司”。\protect\footnotemark}
\end{figure}
\footnotetext{视频画面时间区间：00:17:50--00:17:54。}

\subsection{短板与偏见}

视频没有回避 Linux 当时的明显短板：

\begin{itemize}
\item \textbf{应用太少}：没有像样的 Microsoft Office、Visual Basic 或其他商业软件。视频对这个问题的评论颇有意思：《Byte》在指出这一点之后，紧接着补了一句“Linux 确实包含大量可以用来自己写应用的工具”——这恰好说明了当时 Linux 读者的画像。
\item \textbf{支持困难}：出了问题该查什么文档、该打给谁？这正是 Red Hat、Caldera 这类公司能介入的地方。
\item \textbf{声誉问题}：尽管 Linux 与 UNIX 没有直接血缘，它仍常被与 UNIX 联系在一起，因而继承了 UNIX“晦涩难用”的名声。
\end{itemize}

\begin{warningbox}{管理层的偏见是一种真实的市场摩擦}
视频提到，即便 Linux 增长很快、有不少知名机构采用、性价比也确实明显，公司高管层中仍然存在负面偏见。一位经理可能会问：“你怎么能把关键服务器交给它？”\\
而与此同时，Linux 很可能已经在他公司某台机器上跑着了。一位匿名系统管理员 1998 年对《旧金山纪事报》说的话很能代表这种张力：\\
“很多管理层不信任 Linux，因为它是免费的。但他们让我负责运行服务器，那我就用最有效的。最有效的是 Linux。”
\end{warningbox}

\subsection{本章小结}
Linux 的早期增长由价格优势启动，但真正让它站稳的是可靠性、源代码可访问性和可定制性。它的短板（应用生态、技术支持、声誉）真实存在，也因此给发行版公司和后来的商业生态留下了空间。管理层偏见与技术人员的实际选择之间的张力，预示了 Linux 接下来会以“自下而上”的方式进入企业。

\section{1998：Intel 发现 Linux}

\subsection{来自国家实验室的请求}

1998 年初，Intel 开始收到来自全国各地的奇怪请求。费米国家加速器实验室（Fermi National Accelerator Lab）等机构的工程师和研究人员，要求 Intel 把它的数学库移植到 Linux 上。

在此之前，Intel 与 Linux 几乎没有什么交集，最多修过一两个 bug。这些请求促使 Intel 去核查情况。他们向为小型 ISP 等客户生产白牌服务器的 OEM 打听，得到的反馈是：Linux 的增长正在加速向上。

具体数字很有冲击力：\textbf{大约 20\% 的公共 Web 服务器运行 Linux}。考虑到这个操作系统在桌面上几乎没有任何存在感，这个份额出乎意料地高。

\begin{figure}[H]
\centering
\includegraphics[width=0.72\textwidth]{figures/fig15-oracle-linux.jpg}
\caption{1998 年 7 月 CNET 的短讯：“Oracle offers Linux”。数据库厂商的移植决定，是那个夏天一系列转向的标志性事件之一。\protect\footnotemark}
\end{figure}
\footnotetext{视频画面时间区间：00:25:05--00:25:09。}

这个消息一路传到了安迪·格鲁夫那里。据视频描述，他用相当“丰富多彩”的语言表示，他们应该对此做点什么。

\subsection{看不懂的商业模式}

于是 Intel 高管联系了 Red Hat，在系列会议中试图了解这家公司是怎么运作的。结果他们有点被搞糊涂了：Red Hat 并不拥有其核心软件产品背后的知识产权，而任何人都可以免费获得这份软件。这怎么赚钱？这怎么投资？

Red Hat 的人继续解释。CEO Robert Young 喜欢用番茄酱作比喻：任何有番茄、糖和醋的人都能做番茄酱，但亨氏（Heinz）占据了这个市场 80\% 的份额。言下之意是：产品本身可以复制，但品牌、渠道、支持和服务构成的商业位置不容易复制。

\begin{figure}[H]
\centering
\includegraphics[width=0.42\textwidth]{figures/fig16-future-of-linux.jpg}
\caption{1998 年 7 月 14 日在圣克拉拉举行的“Linux 的未来”座谈会记录。主办方预计到场人数远超预期，约 700--900 人到场。\protect\footnotemark}
\end{figure}
\footnotetext{视频画面时间区间：00:25:16--00:25:21。}

同时，Intel 内部的 Linux 支持者也在分享他们从 ISP 渠道了解到的采用数据。最终，Intel 管理层被说服了，两家公司开始讨论合作。

\subsection{Merced 的焦虑}

Intel 此时还有一个更具体的动机。他们在开发一款叫 Merced 的新芯片，这是其 64 位 IA-64 服务器架构的第一个实现，后来被命名为 Itanium，据说速度极快。

但担忧一直存在：NT 可能会捆住它的全部能力——就像 32 位时代那样，整整耽误了十年。这种担忧还有新的佐证：微软的下一代 NT（当时叫 NT 5.0）从 1998 年推迟到 1999 年，之后又经历更多延期，最终被改名为 Windows 2000。

\begin{importantbox}{Intel 的真实诉求}
Intel 想要的不只是“多一个操作系统能跑在 x86 上”。它想要的是一个\textbf{可以被自己深度参与优化}的操作系统——一个能及时吃掉最新芯片能力、并且不会被某个伙伴卡住的操作系统。Linux 恰好满足这两点：代码开放，且没有一个能对它行使否决权的所有者。
\end{importantbox}

\subsection{本章小结}
1998 年，来自科研机构的移植请求和 ISP 渠道的数据让 Intel 意识到 Linux 在服务器侧的采用已经不可忽视。Red Hat 的“番茄酱”商业模式最终说服了 Intel 管理层，而 Merced/Itanium 的软件依赖焦虑则提供了更直接的动机。Intel 需要 Linux，不只是为了多一条产品线，而是为了摆脱对微软单一路径的依赖。

\section{Linux 之夏与 Intel 的押注}

\subsection{一连串的公开表态}

1998 年夏天，多家大公司接连发布与 Linux 相关的公告。

\textbf{6 月，IBM。} 这家以保守著称的公司宣布，将在其电子商务工具套件中纳入开源 Web 服务器 Apache。Apache 本身不是 Linux，但绝大多数 Apache 实例都跑在 Linux 上。IBM 早在 1996 年就选用了 Apache 为亚特兰大夏季奥运会官网提供支持，这次则是正式背书——开源软件第一次在企业市场获得立足点。

\textbf{7 月 6 日，Oracle 说“不”。} 《InfoWorld》发表了一篇题为“Linux not yet critical mass”的文章，其中 Oracle 连同 Sybase、Informix 都表示没有计划为 Linux 发布数据库软件，理由是没看到客户有强烈需求。

\textbf{7 月 9 日，Oracle 转向。} 仅仅三天后，Oracle 宣布将把 Oracle8 数据库移植到 Linux。导致这次急转弯的重要原因，是 7 月在圣克拉拉举行的一场名为“Linux 的未来”的座谈会。

\textbf{7 月 14 日，座谈会。} 900 多人到场看 Torvalds 等人。会上讨论了开源是否只是一阵风、并敦促用户推动应用厂商移植到 Linux。后来人们把这场座谈会视为一个重要的转折点。

\textbf{9 月，更多厂商跟进。} 到 1998 年 9 月，已经有多家大型公司提供 Linux 移植版本，包括企业系统公司 Computer Associates、软件公司 Corel 和数据库公司 Informix。这个趋势在随后几年持续。

\subsection{Intel 的投资}

Intel 也参加了 7 月那场座谈会。Intel 的 Sunil Saxena 在会上做报告，指出 Linux 正在成为 ISP 的首选操作系统，同时点出了它的几个弱点：可扩展性、高端硬件驱动、大内存支持。他强调 Intel 希望支持 Linux。

这些话很快变成了行动。1998 年 9 月 29 日，Intel 宣布与浏览器公司 Netscape、风投 Benchmark 和 Greylock 一起投资 Red Hat。

\begin{figure}[H]
\centering
\includegraphics[width=0.34\textwidth]{figures/fig17-intel-pressrelease.jpg}
\caption{Intel 1998 年 9 月 29 日的新闻稿，宣布在 ISPCON Fall '98 上推出一系列面向 ISP 的举措，其中包括对 Linux 社区的支持。\protect\footnotemark}
\end{figure}
\footnotetext{视频画面时间区间：00:26:23--00:26:28。}

Intel 的表态包括几层含义：

\begin{itemize}
\item 希望 Intel 芯片运行 Linux 的效果，和运行 Windows NT 及其他系统一样好；
\item 与 Red Hat 合作，提供技术贡献，并为 Linux 社区指派专职的 Intel 工程联络人；
\item 加入 Linux International——一个早期的全球性 Linux 倡导组织。
\end{itemize}

\subsection{一个被低估的多米诺骨牌}

投资公告之后，还有一张关键的多米诺骨牌倒下，这件事在视频里被称为“关键但被低估”。

Benchmark 的一位合伙人同时是戴尔的董事会成员。通过他，Red Hat 得以与 Michael Dell 会面。双方讨论了 Linux 在服务器上的增长、市场格局会如何演变，以及 Linux 有没有可能赢得桌面。

戴尔最终投资了 Red Hat。更重要的是，他们在 1999 年初宣布将在工作站和服务器上预装 Linux——这是第一家这么做的整机大厂。

\begin{figure}[H]
\centering
\includegraphics[width=0.8\textwidth]{figures/fig18-dell-redhat.jpg}
\caption{Red Hat 1999 年 2 月 4 日的新闻稿：戴尔将预装 Red Hat Linux。这是主流系统厂商第一次在整机上预装 Linux。\protect\footnotemark}
\end{figure}
\footnotetext{视频画面时间区间：00:27:08--00:27:12。}

预装的意义远大于投资。投资只是财务行为，预装意味着 Linux 被放到了与 Windows、NetWare 同等的采购选项里。客户买服务器时，开机就是 Linux，不再需要自己折腾安装。

\subsection{微软的反应：从暴怒到“Linux Myths”}

Intel 投资 Red Hat 的消息公开后，科技媒体立刻开始猜测这对 Intel 与微软的关系意味着什么。

Red Hat 高管在《Under the Radar》一书中回忆，微软在 Intel 宣布投资后“went ballistic”（暴跳如雷）。那个周末，Intel 高管接到了多个来自微软对应人员的愤怒电话。

但 1998 年的微软已经没有 1997 年那样的自由度了。1998 年 5 月，美国司法部以垄断行为起诉微软，主要针对浏览器大战期间将 Internet Explorer 与 Windows 捆绑的做法。诉讼和随之而来的审查，迫使微软收敛行为。

微软公开的回应很巧妙：他们发表声明，把 Red Hat 说成是“微软并没有垄断操作系统”的证据。更有意思的是，微软后来还援引 Linux，作为操作系统市场仍有足够竞争、政府不该拆分他们的证据。

但内部并不轻松。1998 年 10 月，一组后来被称为“万圣节备忘录”（Halloween Memos）的微软内部文件泄露，其中把 Linux 这类自由软件的兴起称为对 Windows 的“长期可信威胁”。

\begin{figure}[H]
\centering
\includegraphics[width=0.9\textwidth]{figures/fig19-halloween.jpg}
\caption{“万圣节备忘录”（The Halloween Documents）公开页面。这些内部文件系统分析了开源开发模式的优势，并把 Linux 视为长期威胁。\protect\footnotemark}
\end{figure}
\footnotetext{视频画面时间区间：00:28:29--00:28:33。}

随后微软据称组建了一支“Linux 攻击小组”，试图说服客户 Linux 还没准备好。他们上线了一个叫“Linux Myths”的网页，声称要“从炒作中退后一步，看看现实”，逐条反驳 Linux 信徒们念叨的“神话”，比如“Linux 比 Windows NT 4 性能更好”“Linux 比 Windows NT 4 更可靠”“Linux 是免费的”。

最具争议的一次行动发生在 1999 年 4 月：微软资助第三方实验室 Mindcraft 做基准测试，结果显示 NT 在文件服务器场景下比 Linux 快 2.5 倍，在 Web 服务器场景下快 3.7 倍。

\begin{figure}[H]
\centering
\includegraphics[width=0.84\textwidth]{figures/fig20-mindcraft.jpg}
\caption{Mindcraft 的对比测试报告（左）与测试结果图表（右）。Linux 支持者随后指出该测试的配置并不公平，这一争议持续了很长时间。\protect\footnotemark}
\end{figure}
\footnotetext{视频画面时间区间：00:29:10--00:29:15。}

Linux 支持者随后指出这次测试不公平，尤其是 Linux 一方使用了未经调优的旧版本内核与配置。但无论争议结果如何，这件事本身说明了一点：微软已经把 Linux 当成了需要正面应对的竞争对手。

\subsection{本章小结}
1998 年的“Linux 之夏”由一连串标志性事件组成：IBM 接纳 Apache、Oracle 三天内反转、圣克拉拉的座谈会、Intel 投资 Red Hat、戴尔宣布预装。Intel 的“OS 中立”表态为它自己留出了回旋空间，而微软则从公开否认转向内部紧张，万圣节备忘录和 Mindcraft 测试都是这种紧张的外化。

\section{Port of Choice：瑞士策略与硬件补齐}

\subsection{OS 中立}

面对盖茨的怒火，Intel 的应对方式是宣布自己“OS 中立”。Linux、Novell 的 NetWare、Sun 的 Solaris、微软的 Windows 2000，Intel 希望它们都能在自己的硬件上跑得好。

Intel 甚至与另外二十家公司一起，试图用一种叫 Uniform Driver Initiative 的机制统一当时三十来种 UNIX 变体的驱动模型，方便厂商为 Intel 硬件写驱动。这件事最终没有成功。

这种策略后来在 2010 年被正式命名为“Port of Choice”，但 Intel CTO Pat Gelsinger 在 2001 年就说这是公司长期坚持的策略。最终操作系统厂商不得不接受这个现实，因为当时的 Intel 芯片没有真正的替代品。

\begin{knowledgebox}{什么叫“Port of Choice”}
字面意思是“移植的首选目标”：任何操作系统如果想获得最大的硬件覆盖面，就必须优先支持 Intel 架构。这个位置的价值在于，它让 Intel 不再需要依附于某一个操作系统——相反，操作系统需要来适配 Intel。
\end{knowledgebox}

\subsection{补齐硬件：从 4 路到 8 路}

软件之外，硬件端也在补齐服务器所需的能力。

1998 年中，Intel 推出了 Pentium Pro 的后续产品 Pentium II Xeon，主攻中端工作站与服务器市场，特点是更高的主频和 512 KB 的 L2 缓存。

几个月后的 1999 年 3 月，Intel 发布 Pentium III Xeon。相比前代，它的速度提升不大，但支持通过名为 Profusion 的芯片组构建 8 路系统。这意味着厂商可以在一个机箱里对称地放八颗 CPU——突破了此前 4 路系统的天花板，而那个天花板正是当时数据中心性能问题的来源之一。

\begin{figure}[H]
\centering
\includegraphics[width=0.72\textwidth]{figures/fig21-xeon-8way.jpg}
\caption{8 路 Xeon 系统的结构示意：八颗处理器通过 Profusion 芯片组共享内存与 I/O 总线。\protect\footnotemark}
\end{figure}
\footnotetext{视频画面时间区间：00:30:57--00:31:02。}

这里有一个关键细节：8 路能力需要操作系统配合才能发挥作用。Windows NT/2000 支持它，而 Linux 当时还有问题。这解释了为什么 Intel 同时要在两个操作系统上下注——硬件能力已经就位，软件能不能跟上决定了它能不能卖出去。

\subsection{高端仍有护城河}

高端 UNIX 服务器仍然保留着关键优势。最突出的是内存寻址能力：大型数据库、数据仓库、CAD 系统需要的主内存远超过 32 位 x86 架构所能支持的上限。

但凭借 Pentium II Xeon 和 Pentium III Xeon 这两代产品，中低端 UNIX 服务器的侧翼已经完全暴露在竞争之下。

\subsection{本章小结}
Intel 用“OS 中立”的瑞士策略避免被任何一家操作系统绑定，同时通过 Xeon 系列和 Profusion 把 x86 服务器的多路扩展能力从 4 路推到 8 路。硬件补齐之后，剩下的瓶颈就落在软件侧：NT 的 SMP 上限和 Linux 的 SMP 缺陷，成了下一阶段竞争的主战场。

\section{Web 让天平倾斜：scale-out 的经济学}

\subsection{Sun 的误判}

Sun 最初是欢迎 Linux 革命的。这与他们“开放”的口号一致：他们把 Java 开发工具包移植到 Linux，还发布了可运行 Linux 的工作站硬件。他们给人的总体印象是：Sun 不在乎一台机器跑 Linux 还是 Solaris，只要它是 SPARC 机器就行。

但事实证明这是一个误判：它不必是 SPARC 机器。

提供网站服务、邮件托管、防火墙这类工作，与运行数据库或分析软件不是一回事。这类负载可以水平切分到很多台小机器上。企业逐渐意识到：用一批 Intel 芯片加 Linux 的节点横向扩展（scale out），在财务上比纵向扩展（scale up）一台大型 UNIX 服务器更划算。

\subsection{泡沫破裂加速了这一趋势}

互联网泡沫破裂后，这个趋势突然加速。风险投资枯竭，创业公司开始倒闭。节俭成了主旋律，企业通过购买倒闭公司清算出来的二手设备，或者直接改用 Linux 来省钱。

亚马逊在 2001 年做了一次大规模的 Linux 迁移，并在季度 SEC 文件中提到，这估计每季度为他们节省了 1700 万美元。当被问及具体做了什么时，亚马逊发言人的回答很有意思：

\begin{figure}[H]
\centering
\includegraphics[width=0.62\textwidth]{figures/fig24-amazon-linux.jpg}
\caption{2001 年 Yahoo News 转载的路透社报道：“Amazon Says Saving Money with Shift to Linux”。\protect\footnotemark}
\end{figure}
\footnotetext{视频画面时间区间：00:32:46--00:32:58。}

\begin{quote}
“We've always been pretty closed-mouthed about technology on the back-end stuff, partly because it's pretty steep tuition and we don't want other people going to school on our tuition.”
\end{quote}

翻译过来是：我们对后端技术一向守口如瓶，部分原因是学费很贵，我们不想让别人拿我们的学费去上学。

\subsection{E*Trade 的账本}

在线券商 E*Trade 在 2006 年公布了一份很有说服力的案例研究。它对比了两套方案：

\begin{table}[H]
\centering
\caption{E*Trade 案例中的两套方案对比}
\begin{tabular}{lccc}
\toprule
方案 & 硬件成本 & 可服务用户数 & 每用户成本 \\
\midrule
Sun Enterprise 4500 & \$250,000 & 300--400 & 约 \$714 \\
Linux + Intel 单机 & \$4,000 & 180 & 约 \$22 \\
Linux + Intel 两台 & \$8,000 & 360 & 约 \$22 \\
\bottomrule
\end{tabular}
\end{table}

单台 Linux 机器只能服务 180 人，看起来不如 Sun 的 300--400 人。但解决办法很简单：买两台。总价 8000 美元，服务 360 人，仍然只有 Sun 方案的 3.2\%。

\begin{figure}[H]
\centering
\includegraphics[width=0.92\textwidth]{figures/fig-gen-cost-per-user.pdf}
\caption{按每用户硬件成本计算，scale-out 方案的差距约为 32 倍。}
\end{figure}

\begin{figure}[H]
\centering
\includegraphics[width=0.74\textwidth]{figures/fig23-etrade.jpg}
\caption{E*Trade 2006 年的案例研究与 Sun Enterprise 4500 服务器。文中提到这次迁移每年节省约 1300 万美元。\protect\footnotemark}
\end{figure}
\footnotetext{视频画面时间区间：00:33:09--00:33:16。}

从 2001 年开始，搭载 Intel 芯片的 Linux 服务器逐渐有能力稳定替换 Sun 的设备，覆盖从机票比价到大宗商品合约交易再到金融计算的各类场景。

\subsection{Sun 的应对与失败}

Sun 也尝试过自己的横向扩展产品线 Netra。2001 年的 Netra X1 售价低至 1000 美元，但通用 x86 硬件仍然提供了更好的每美元性能。

视频的判断相当直接：这家公司始终没有针对“从 SPARC 上的 Solaris 转向 Intel 上的 Linux”这一转变，形成过像样的回应。

\begin{figure}[H]
\centering
\includegraphics[width=0.8\textwidth]{figures/fig22-serveros-share.jpg}
\caption{1995--2000 年各类服务器操作系统的全球出货量份额。Linux 从 1996 年几乎为零增长到 2000 年的约四分之一，而 UNIX 的份额持续收缩。\protect\footnotemark}
\end{figure}
\footnotetext{视频画面时间区间：00:32:21--00:32:27。}

\subsection{本章小结}
真正的转折发生在 Web 负载上。这类负载可以水平切分，使得“多台廉价机器”在财务上碾压“一台大型 UNIX 服务器”。互联网泡沫破裂又加速了节俭采购的趋势。E*Trade 与亚马逊的案例说明，这种替换的驱动力来自单位成本的数量级差异，而非单机性能。

\section{向上攀升：Linux 2.4、2.6 与企业级}

低端市场拿下之后，Linux 开始向上进攻。

\subsection{IBM 的十亿美元}

2000 年底，IBM CEO 郭士纳（Lou Gerstner）认为 Linux 会像互联网一样改变世界，宣布了一项大规模的 Linux 计划：投入 10 亿美元，投入 1500 名员工。

\subsection{Linux 2.4：补上 SMP 这一课}

2001 年 1 月，Linux 2.4 发布。Red Hat 很快在后续版本中采用。它比预期晚了一年，但在缩小与高端 UNIX 服务器和 NT 的若干性能差距方面，确实兑现了承诺。

最关键的一项改进是多处理器支持。1999 年 Mindcraft 测试中的一个主要争议点，就是 Linux 对多处理器机器的支持。那次测试（以及后续测试）显示：Linux 2.2 在单系统内超过两颗 CPU 时扩展性很差。而 Windows NT 4.0 和 2000 分别支持 4 路、8 路乃至 32 路。缺少这个能力是一个重大弱点。

Torvalds 和社区因此着手改进。2.4 版处理了自旋锁（spin lock）等问题，使 Linux 能够在某种程度上有效使用 16 甚至 32 颗处理器——不过视频作者也坦率地说，他不确定在那么高的核数下实际跑得怎么样。

这个改进并没有让 Linux 一夜之间挑战高端 UNIX 机器，但至少具备 4 路或 8 路的能力，为“Linux on Intel”（也就是 Lintel）进入中端计算任务扫清了道路。

\begin{figure}[H]
\centering
\includegraphics[width=0.68\textwidth]{figures/fig25-linux24.jpg}
\caption{2001 年 1 月 Linus Torvalds 在 Linux 内核邮件列表上发布 2.4.0 的邮件。通篇是自嘲式的幽默，最后一句是“别在这几天报 bug，反正我也不在乎”。\protect\footnotemark}
\end{figure}
\footnotetext{视频画面时间区间：00:34:26--00:34:32。}

\subsection{经济学变得极具说服力}

Lintel 的经济账忽然变得非常吸引人。一台戴尔服务器，四颗 Xeon，跑 Red Hat Linux，售价 24000 美元；价格是同等 Sun Fire UNIX 机器的一半，性能却一样好。

\begin{knowledgebox}{为什么“一半价格、同等性能”是毁灭性的}
服务器采购不是一次性的。企业还要付电费、机架空间、运维人力、软件许可。价格减半意味着同样的预算能部署两倍节点，或者同样的负载能省下一半资本开支。当这个差距持续存在时，采购决策几乎不需要讨论——这也是为什么 Sun 的应对空间越来越小。
\end{knowledgebox}

Red Hat 和 IBM 立即开始瞄准 Sun 的 Solaris 客户群。惠普和 SGI 则开始给使用其专有硬件的客户提供迁移到 Linux 的选项——理由很务实：与其让客户流失，不如保住客户关系。

\subsection{2002 年的市场数据}

2002 年对服务器市场来说是艰难的一年，但 Linux 机器仍然在增长，而且是以 RISC 系统为代价。

Gartner 的数据显示：2002 年 Linux 服务器硬件收入增长 63\%，而 UNIX 硬件收入下降 11\%。

从绝对金额看，UNIX 硬件仍然更大——2002 年为 170 亿美元，Linux 为 20 亿美元——但这只是因为 UNIX 是高端机器，单价更高。压力已经非常明显：到 2004 年，Sun 服务器的平均售价从 22920 美元跌到了 15860 美元，跌幅 31\%。

\subsection{Linux 2.6：第一个“企业就绪”的内核}

2003 年底，Linux 2.6 发布。视频称它为一个非常标志性的内核：速度、网络、对最多 64 颗处理器的支持等各方面都有改进，可以视为第一个“企业就绪”的 Linux 内核。

\begin{figure}[H]
\centering
\includegraphics[width=0.8\textwidth]{figures/fig27-linux26.jpg}
\caption{The Register 对 Linux 2.6 发布的报道，“Penguins are go”。\protect\footnotemark}
\end{figure}
\footnotetext{视频画面时间区间：00:38:53--00:38:58。}

更需要强调的是时机：它恰好赶上了一批 Windows NT 4.0 用户开始升级周期。这些用户正处在“单身待嫁”的状态，而 IBM、惠普、戴尔等厂商在拼命把客户往 Linux 上引。

IBM 尤其积极：培训了数千人，帮助开发者把应用重写到 Linux 上，还向潜在迁移者提供大笔营销预算和折扣机器。

\subsection{本章小结}
Linux 2.4 补上了 SMP 这一课，让 Lintel 得以进入中端市场；价格优势加上性能追平，使得“一半价格、同等性能”成为常态；Gartner 的 2002 年数据表明 Linux 的增长直接来自 UNIX 的份额流失。Linux 2.6 则把 Linux 从“能用”推向“企业就绪”，并且恰好卡在 NT 4.0 的升级窗口上。

\section{64 位：最后一堵墙倒下}

\subsection{Trillian 与 Itanium}

Intel 很早就希望 Linux 支持 64 位，也就是 IA-64 或 Itanium。1999 年 4 月，Intel 启动了 Project Trillian——一个由多家公司组成的联盟，成员可以保密地获得 Itanium 原型，以便开发 Linux 移植版本。2000 年 2 月，该项目提前向社区发布了 IA-64 Linux 源代码。

视频称 Trillian 是开源软件开发史上的一个里程碑。但 Itanium 最终并没有取代 32 位 x86，尽管惠普确实采用了它来替换其专有 UNIX 变体所使用的处理器。

\subsection{AMD 完成了这次跃迁}

真正开启 64 位过渡的是 AMD，而不是 Intel。2003 年中，AMD 推出了开创性的 Opteron 系列芯片。2004 年，Intel 发布了 Nocona Xeon，采用了 AMD 的向后兼容 64 位扩展。

\begin{figure}[H]
\centering
\includegraphics[width=0.82\textwidth]{figures/fig26-opteron.jpg}
\caption{AMD Opteron 处理器。它用向后兼容的方式把 64 位引入 x86 生态，成为压垮 UNIX 最后一道架构壁垒的产品。\protect\footnotemark}
\end{figure}
\footnotetext{视频画面时间区间：00:37:38--00:37:43。}

\subsection{为什么这是“最后一堵墙”}

64 位一直是高端 UNIX 系统相对 x86 的最后优势之一。原因很具体：

\begin{itemize}
\item 32 位 x86 计算机最多寻址约 4 GB 内存；
\item 而数据库、科学模拟这类高端应用需要一次性驻留在内存里的数据量，远超这个上限；
\item 转向 64 位后，x86 系统可以寻址数十甚至数百 GB 内存。
\end{itemize}

\begin{figure}[H]
\centering
\begin{tikzpicture}[font=\small]
\draw[fill=blue!6, draw=blue!60!black] (0,0) rectangle (6.4,1.1);
\node at (3.2,0.55) {32 位地址空间：$2^{32}$ B = 4 GB};
\draw[fill=green!8, draw=green!50!black] (0,1.6) rectangle (6.4,2.7);
\node at (3.2,2.15) {64 位地址空间：$2^{64}$ B（实际受硬件限制，可达数十 TB）};
\node[anchor=east] at (-0.15,0.55) {x86};
\node[anchor=east] at (-0.15,2.15) {x86-64};
\end{tikzpicture}
\caption{地址空间是高端应用的一道硬性门槛。把 4 GB 的墙拆掉之后，UNIX 在架构层面最后的结构性优势也随之消失。}
\end{figure}

有人最初认为 Intel 采纳 64 位的方式是一种权宜之计，只是通向 Itanium 的过渡。但结果无可否认：x86-64 砸碎了围绕 UNIX 的最后几堵墙。视频用了一个很好懂的说法：Lintel 的“蛮族”已经进入城堡。

\subsection{本章小结}
64 位寻址能力是高端 UNIX 系统最后的结构性优势。Intel 主导的 IA-64 路线（Itanium）未能成为主流，反而是 AMD 用向后兼容的方式把 64 位带进 x86。当 x86-64 与 Linux 结合之后，UNIX 在架构层面已经没有不可替代的能力。

\section{Windows 的反击：Active Directory 的锁定}

\subsection{微软的营销手段}

面对 Linux 的攻势，微软的第一反应是打营销战。它的核心论点是：大家已经在用 Office，与其折腾迁移，不如直接买 Windows Server 许可证。

2004 年，微软发起了一场叫“Get the facts”的营销活动，发布声称 Windows 优于 Linux 的研究。这些研究很快引来争议，被指不够“中立”，甚至可能误导。

\subsection{真正有效的武器}

视频作者认为，真正有效的做法不是营销，而是一款叫 Active Directory（AD）的产品带来的深度锁定。Stratechery 的 Ben Thompson 后来把它称为“有史以来最有价值的免费产品之一”。

AD 随 Windows 2000 一起发布。它的作用是认证和授权一个 Windows 网络里的所有用户和资产：可以推送更新、按照组织架构给用户组分配安全权限、验证身份。

具体到日常运维，AD 带来的变化是：

\begin{itemize}
\item 管理员可以为一个新上任的经理开通账户、验证凭据、授予对应工具权限；
\item 用户用一套凭据就能登录各种系统；
\item 一堆零散的 Windows 机器被整合成一个可统一管理的组织。
\end{itemize}

\begin{figure}[H]
\centering
\includegraphics[width=0.6\textwidth]{figures/fig28-activedirectory.jpg}
\caption{Windows 2000 中的 Active Directory 用户和计算机管理界面。\protect\footnotemark}
\end{figure}
\footnotetext{视频画面时间区间：00:40:04--00:40:08。}

关键点在于：AD 让微软掌握了“身份”这一层。身份深深嵌入组织内部，因此把切换成本抬高到了“换个操作系统”之外——你要换的不只是服务器系统，还有整个组织的账号体系、权限模型和运维流程。

\begin{figure}[H]
\centering
\includegraphics[width=0.82\textwidth]{figures/fig29-gartner-decline.jpg}
\caption{2011--2017 年 UNIX 服务器安装量与占比的持续下滑。即便在被 Linux 击败十年之后，UNIX 也没有完全消失，但下降趋势从未逆转。\protect\footnotemark}
\end{figure}
\footnotetext{视频画面时间区间：00:41:41--00:41:46。}

\subsection{结果：谁也没有吃掉谁}

所以在那些预测 Linux 会推翻微软的人面前，Windows Server 的销售继续增长。2005 年，Windows 的服务器收入首次追平 UNIX，次年超过。

Linux 也在继续增长，但它的增长主要不是来自吃 Windows，而是来自吃 UNIX，以及在 Windows 从未立足的新市场里扩张：超级计算、高性能计算、Web，以及 2010 年代的云计算。

\subsection{本章小结}
微软真正有效的防御不是营销战，而是 Active Directory 建立的深度锁定。它把竞争层次从“操作系统替换”抬高到“组织身份体系重建”，从而保住了 Windows 在企业市场的地位。最终格局是：Linux 吃掉了 UNIX 的大部分市场，Windows 保住了自己的地盘，而 UNIX 在缓慢衰退中退守少数领域。

\section{总结与延伸}

\subsection{作者的结论}

UNIX 服务器的销售在随后几年缓慢下滑，但并没有完全消亡。即便在今天，仍有客户在使用 Solaris、HP-UX 或 AIX。这些滞后应用大多集中在银行和医疗行业，原因各异——对它们来说，迁移到 Linux 要么太难，要么风险太高。

那么 UNIX 为什么会输？视频作者给出的答案很明确：

\begin{quote}
“I might propose that it had less to do with technology than business model. These vertically integrated kingdoms worked fine years ago until Intel arrived with its high volume economics.”\\
（我认为这与技术关系不大，更多是商业模式的问题。这些垂直整合的王国在多年前运转良好，直到 Intel 带着它的高销量经济学到来。）
\end{quote}

作者给出的顺序是：先是 Windows NT 重创了 UNIX，然后 Linux 打得更狠——因为它能给客户 UNIX 拥有的一切，却没有锁定。一旦 Linux 变得足够好，UNIX 厂商自己也不得不转向 Linux。

\subsection{把机制压缩成三层}

把整条线索压缩，这场替代实际上同时发生在三个层面上：

\begin{importantbox}{三层替代机制}
\textbf{成本层。} Intel 把 PC 的高销量飞轮带进服务器，使得同样性能的硬件价格下降一个数量级。SHV 参考设计进一步把整机成本压到 OEM 可以直接组装的水平。\\
\textbf{生态层。} GPL 让 Linux 成为一个可以被所有公司共同投资、却不能被任何一家公司私有化的平台。这让 IBM、Intel、Red Hat 等利益并不一致的公司，能够在一个方向上共同投入。\\
\textbf{锁定层。} UNIX 的锁定来自专有软硬件的捆绑；Lintel 的锁定则几乎为零——操作系统可以换、硬件可以换、厂商可以换。低锁定本身就是一种产品特性，它让采购决策变成纯粹的成本性能计算。\\
\end{importantbox}

这三层是互相加强的：成本下降带来更多装机量，装机量吸引更多软件支持，软件支持又进一步降低成本。而 UNIX 厂商的垂直整合结构，恰恰无法参与这个循环——它们的收入模型建立在“高单价、高锁定”之上，任何降价都会直接侵蚀利润。

\subsection{一个反直觉的结果}

值得注意的是，Linux 的胜利并不等于 UNIX 客户的“解放”以另一种锁定告终。视频在结尾指出，Linux 的增长后来主要发生在 Windows 从未立足的新市场——超级计算、HPC、Web 和云计算。

而在企业市场内部，微软用 Active Directory 建立了新的锁定层，把身份体系变成了切换成本的主要来源。于是最终格局是三分天下：Linux 拿下 Web 与新兴负载，Windows 守住企业身份与桌面协同，UNIX 退守少数监管严格或迁移风险极高的行业。

\subsection{可以带走的判断}

\begin{enumerate}
\item \textbf{技术指标的领先是必要条件，不是充分条件。} Pentium Pro 在浮点上始终落后于最好的 RISC，Linux 2.2 的 SMP 也很糟糕。它们赢，是因为“够用 + 便宜 + 不被锁定”的组合。
\item \textbf{治理结构是产品竞争力的一部分。} GPL 不是法律细节，它决定了 Linux 能否吸收大公司的贡献而不被碎片化。
\item \textbf{锁定层的迁移比技术替代更慢，也更关键。} UNIX 的技术劣势在 1998 年就很明显，但它真正退出主流用了十多年；同样，Linux 在服务器上占了上风，也没能撼动 Active Directory 建立的锁定。
\item \textbf{伙伴的动机需要单独分析。} Intel 支持 Linux，根源在于它对微软的结构性依赖；理解这一点，比记住任何一次发布会都更有解释力。
\end{enumerate}

\subsection{开放的问题}

视频没有展开、但值得继续追问的几个方向：

\begin{itemize}
\item 如果没有 AMD 的 Opteron，x86 的 64 位迁移会晚多久？Intel 的 Itanium 路线是否可能成功？
\item 如果 GPL 换成一个更宽松的许可证（如 BSD），Linux 是否会重演 UNIX 的碎片化命运？
\item 云计算时代的“锁定层”是否已经从操作系统转移到了 API 与数据格式上？这与 Active Directory 的故事是否同构？
\end{itemize}

\subsection{参考}

\begin{enumerate}
\item Asianometry. \emph{The Lintel Conquest of UNIX}. YouTube, 2026. \url{https://www.youtube.com/watch?v=ksx5g5WvAIw}
\item Young, Robert, and Wendy Goldman Rohm. \emph{Under the Radar: How Red Hat Changed the Software Business---and Took Microsoft by Surprise}. Coriolis Group, 1999.
\item Colwell, Bob. \emph{The Pentium Chronicles: The People, Passion, and Politics Behind Intel's Landmark Chips}. Wiley-IEEE Computer Society Press, 2006.
\item Raymond, Eric S., ed. \emph{The Halloween Documents}. \url{https://www.catb.org/~esr/halloween/}
\item United States District Court for the District of Columbia. \emph{Plaintiffs' Joint Proposed Findings of Fact, United States v. Microsoft Corp.}, 1999.
\item Thompson, Ben. ``Microsoft and Software Survival.'' \emph{Stratechery}. \url{https://stratechery.com/2020/microsoft-and-software-survival/}
\item Slater, Michael. ``Intel Boosts Pentium Pro to 200 MHz: Integer Performance Beats the Best RISCs.'' \emph{Microprocessor Report}, vol. 9, no. 15, 13 Nov. 1995.
\item Schlender, Brent. ``Intel's \$10 Billion Gamble.'' \emph{Fortune}, vol. 146, no. 9, 11 Nov. 2002.
\item Takahashi, Nobuo, and Tomofumi Takamatsu. ``UNIX License Makes Linux the Last Major Piece of the Puzzle.'' \emph{Annals of Business Administrative Science}, vol. 12, 2013, pp. 123--137.
\end{enumerate}

%% --- 正文内容结束 --- %%

\end{document}
